% Common lab-report packages and font shapes (Computer Modern, OT1).
\documentclass[12pt]{article}
\usepackage[letterpaper,margin=1in]{geometry}
\usepackage[utf8]{inputenc}
\usepackage[english]{babel}
\usepackage{amsmath,amssymb,amsthm,mathtools}
\usepackage{gensymb}
\usepackage{upgreek}
\usepackage{graphicx}
\usepackage{float}
\usepackage{wrapfig}
\usepackage{caption}
\usepackage{subcaption}
\usepackage{booktabs,multirow,tabularx,array,longtable,makecell}
\usepackage{xcolor}
\usepackage{enumitem}
\usepackage{fancyhdr}
\usepackage{titlesec}
\usepackage{setspace}
\usepackage{siunitx}
\usepackage[version=4]{mhchem}
\usepackage{chemfig}
\usepackage{textcomp}
\usepackage{url}
\usepackage[numbers,sort&compress]{natbib}
\usepackage{hyperref}
\usepackage{cleveref}
\pagestyle{fancy}
\fancyhf{}
\rhead{Lab report}
\lfoot{\thepage}
\theoremstyle{definition}
\newtheorem{definition}{Definition}
\begin{document}
\section{Shapes and symbols}\label{sec:a}
Roman, \textit{italic}, \textbf{bold}, \textbf{\textit{bold italic}}, \textsc{Small Caps},
\textsf{sans \textbf{bold} \textit{italic}}, \texttt{typewriter}, \emph{emph}, \underline{under}.
{\small small} {\footnotesize footnote} {\scriptsize script} {\large large} {\Large Large}
{\LARGE LARGE} {\huge huge} {\Huge Huge}.
Degrees: 25\,\degree C, \textdegree, \SI{25}{\celsius}, \si{\micro\litre}, $\upmu$, \ang{30}.
Maths: $\alpha + \beta \leq \sum_{i=1}^{n} x_i^2$, $\mathbb{R}$, $\mathcal{L}$,
$\mathfrak{g}$, $\boldsymbol{v}$, $\vec{F} = m\vec{a}$, $\Delta H^\circ$.
\begin{align}
  K_\text{eq} &= \frac{[\ce{C}][\ce{D}]}{[\ce{A}][\ce{B}]} \label{eq:k}\\
  \Delta G &= -RT\ln K
\end{align}
\begin{definition}A limiting reagent is consumed first.\end{definition}
See \cref{eq:k}, \cref{tab:a} and \cref{fig:a}\footnote{A footnote.}. \url{https://example.org}
\begin{itemize}[noitemsep]\item one \item two\end{itemize}
\begin{enumerate}[label=(\alph*)]\item a \item b\end{enumerate}
{\color{red}red} {\color{blue}blue} \colorbox{yellow}{box} \textcolor{gray}{gray}
\ce{^{1}H} NMR \ce{CDCl3}: $\delta$ 7.26. \ce{Cu^2+ + 2e- <=> Cu}. \ce{NaCl(aq)}.
\begin{table}[H]\centering\caption{Table}\label{tab:a}
\begin{tabularx}{\linewidth}{lXS[table-format=2.1]}\toprule
\multirow{2}{*}{A} & \makecell{multi\\line} & 1.0 \\ & x & 22.5\\\bottomrule
\end{tabularx}\end{table}
\begin{figure}[H]\centering
\begin{subfigure}{0.45\linewidth}\centering\chemfig{H-C(-[2]H)(-[6]H)-OH}\caption{Methanol}\end{subfigure}
\begin{subfigure}{0.45\linewidth}\centering\chemfig{*6(-=-(-NO_2)=-=)}\caption{Nitrobenzene}\end{subfigure}
\caption{Structures}\label{fig:a}\end{figure}
\schemestart A \arrow{->[\ce{H+}]} B \schemestop
\onehalfspacing Text. \cite{key}
\begin{thebibliography}{9}\bibitem{key} A. Author, \emph{Book}, 2020.\end{thebibliography}
\end{document}
